\documentclass[../../main.tex]{subfiles}
\begin{document}

\begin{flushright}
  \footnotesize\textit{【自動文字起こし・要確認】}
\end{flushright}

\noindent
{\bf [解答]} 点 $X$ の位置ベクトルを $\vec{x}$ で定める．$\vec{d} = \overrightarrow{PA} + \overrightarrow{PB} + \overrightarrow{PC}$ とし，$G$ を $\triangle ABC$ の重心とすると
\begin{align}
  \vec{d} &= \vec{a} + \vec{b} + \vec{c} - 3\vec{p} \\
  &= 3(\vec{g} - \vec{p})
\end{align}
\begin{align}
  \therefore |\vec{d}| = 3|\vec{g} - \vec{p}| = 3|\overrightarrow{PG}|
\end{align}
したがって，定円の中心 $O$，直線 $OG$ と定円の交点のうち，$G$ と $O$ に関して反対側のものを $E$ として，$P = E$ の時 与式は最大．

\newpage

\end{document}
